\section{State of the Art}
\label{sec:state_of_the_art}

This section compares previous studies according to their optimization target, the gripper properties treated as fixed, and their relevance to
task-specific gripper configuration. The reviewed literature can be divided into four groups: grasp planning with fixed end-effectors, multi-suction
grasping, two-finger gripper design, and comparative studies of multi-objective optimization algorithms.

\subsection{Grasp Planning with Fixed End Effectors}

Most robotic grasping research treats the end-effector geometry as a fixed input and searches for a feasible grasp pose on the workpiece. Zurbr\"ugg
et al.\ \cite{Zurbruegg2025GraspQP}, for example, optimize grasp selection while assuming that the geometry and contact configuration of the gripper
have already been defined. Analytical grasp-planning methods evaluate candidate contacts using geometric, force-closure, or form-closure conditions,
whereas data-driven methods estimate grasp quality from training data. Despite their methodological differences, both approaches primarily vary the
grasp pose rather than the gripper geometry.

Although Todescato et al.\ \cite{Todescato} establish gripper geometry as an optimization variable, their formulation is limited to a single suction
cup. Their objective-function evaluation is methodologically similar to the approach used in this thesis, as suction-contact feasibility is assessed
using quasi-static criteria related to the methods presented by Kleeberger et al.\ \cite{Kleeberger} and the compliant suction-contact model of
Mahler et al.\ \cite{SC_quasi_static}. However, their formulation primarily considers the availability of valid suction grasps and does not explicitly
evaluate the distribution/coverage of the resulting grasp set. The present thesis extends this approach to a multi-suction-cup layout by optimizing
the spatial positions of the individual cups and evaluating valid grasps using orientation and spatial diversity metrics. These metrics distinguish
between layouts that provide a similar number of valid grasps but differ in how broadly those grasps are distributed across positions and
orientations.

The studies therefore address complementary but different questions. Grasp planning determines how an existing gripper should interact with a
workpiece, whereas gripper-configuration optimization determines which geometry should be available before grasp planning is performed. Improving the
grasp planner cannot fully compensate for a suction layout that cannot establish sufficient seals or for a jaw contour that cannot reach suitable
contact regions. Consequently, fixed-gripper grasp-planning results do not directly resolve the end-effector design problem considered in this thesis.

\subsection{Comparison of Multi-Suction Grasping Studies}

Existing studies on multi-suction grippers address grasp prediction, experimental stability assessment, and load-aware motion planning, but generally
assume predefined suction-cup layouts. Schillinger et al.\ \cite{Schillinger2023Model-Free} predict grasp quality from RGB-D images and optimize the
grasp pose and gripper selection from a predefined set of multi-suction grippers. Their method selects among existing layouts rather than optimizing
the positions of the individual cups.

Geng et al.\ \cite{Geng2025Experimental} experimentally investigate how suction-cup size, geometric arrangement, centre of gravity offset, and robot
acceleration influence grasp stability. Although different arrangements are evaluated, these are predefined experimental configurations rather than
layouts generated through multi-objective optimization.

Lee et al.\ \cite{Lee2024Grasp} model the load distribution among multiple suction cups using a quadratic program. They incorporate the resulting
grasp-failure constraints into time-optimal trajectory planning to support the fast and reliable manipulation of heavy objects. Their work focuses on
load prediction and motion planning for an existing gripper configuration, not on optimizing the cup arrangement.

These studies therefore address complementary aspects of multi-suction grasping. Schillinger et al.\ focus on grasp-pose and gripper selection, Geng
et al.\ examine the effects of gripper configuration and handling conditions experimentally, and Lee et al.\ model load distribution during dynamic
manipulation. However, none of these studies formulates the positions of the individual suction cups as decision variables in a multi-objective
layout-optimization problem.

\subsection{Comparison of Two-Finger Gripper Studies}

Research on two-finger gripper design includes the optimization of compliant fingers, underactuated mechanisms, and task-specific jaw geometries. Sun
et al.\ \cite{Sun2022LARG} use three-dimensional topology optimization to design compliant continuum fingers and a lightweight gripper base. The
resulting fingers deform passively to adapt to objects with different shapes. Similarly, Xie et al.\ \cite{Xie2024Design} optimize the topology of
flexible fingers while accounting for contact behaviour and contact-stress constraints. Their objective is to achieve adaptive, damage-free grasping
of fragile and geometrically variable objects. Both studies focus on structural compliance and deformation, whereas this thesis considers the
geometric optimization of rigid jaw contours.

Dong et al.\ \cite{Dong2018Geometric} optimize the geometry and transmission system of an underactuated tendon-driven gripper. Their design variables
include phalange dimensions, palm width, joint-mandrel radii, and tendon-routing parameters. The optimization aims to improve force transmission and
grasp stability. Although this approach optimizes gripper geometry, it concerns the dimensions and actuation mechanism of an articulated gripper
rather than the contact contour of a rigid parallel jaw.

Reichert \cite{reichert2024greiferoptimierung} developed the direct predecessor of the parallel-jaw optimization framework continued in this thesis. The work
represents the contact contour of a rigid parallel jaw using a parametric spline and evaluates candidate contours through a collision-map-based
procedure. Intersections close to the intended workpiece boundary contribute positively to a geometric key performance indicator (KPI), whereas
intersections extending into the workpiece interior are penalized. The KPI therefore measures the geometric conformity of the generated jaw contour
at a specified evaluation configuration rather than providing a complete force- or wrench-based assessment of grasp stability.

Reichert formulates the jaw-design problem using two fixed workpiece orientations, each contributing an objective to the optimization. This
demonstrates the feasibility of combining spline-based jaw generation, geometric KPI evaluation, and multi-objective black-box optimization. However,
the formulation is restricted to the two investigated evaluation configurations.

This thesis continues Reichert's work by extending the evaluation from two fixed configurations to a set of \(n\) predefined grasp points. The grasp
points remain fixed and are not optimization variables; the decision variables are the parameters defining the spline-based jaw contour. Each
candidate contour is evaluated at every predefined grasp point, producing one geometric KPI for each point. All \(n\) KPIs are retained as separate
objectives and maximized simultaneously, resulting in an \(n\)-objective optimization problem with an \(n\)-dimensional objective space.

The optimization therefore seeks Pareto-optimal rigid jaw contours that provide different compromises in geometric conformity across the predefined
grasp points. Compared with the studies on compliant and underactuated grippers, the formulation concerns rigid task-specific contact contours.
Compared with Reichert's direct predecessor formulation, the principal extension is the evaluation and optimization of one separate objective for
each of the \(n\) predefined grasp points.

\subsection{Comparative Studies of Multi-Objective Optimization Algorithms}
\label{subsec:comparative_optimization_studies}

Comparative studies show that optimizer performance depends strongly on the problem and experimental setup. Tanabe and Ishibuchi \cite{Tanabe2020An}
evaluate evolutionary algorithms on a suite of real-world multi-objective problems with varied characteristics, while Tian et al.\
\cite{Tian2019Automated} propose selecting algorithms based on problem features. Ishibuchi et al.\ \cite{Ishibuchi2022Difficulties} further
demonstrate that termination conditions, population sizes, performance indicators, and test problems can change algorithm rankings. Audet et al.\
\cite{Audet2018Performance} emphasize that indicators assess different properties of Pareto-front approximations, such as convergence, distribution,
and spread.

For expensive optimization, Qin et al.\ \cite{Qin2019Bayesian} compare two combinations of evolutionary and Bayesian optimization on ZDT and DTLZ
problems and find that Bayesian evolutionary optimization is more efficient under their experimental conditions. Santoni et al.\ \cite{Santoni2023Comparison}
compare high-dimensional Bayesian optimization methods on single-objective BBOB functions and report that performance varies with dimensionality,
function landscape, and evaluation budget. These findings are therefore specific to the investigated problems and settings.

Such results cannot be transferred directly to the collision-aware suction-layout and spline-based jaw-contour problems considered in this thesis.
Therefore, this thesis compares selected multi-objective evolutionary and Bayesian algorithms under common evaluation budgets, initialization
conditions, repeated runs, and statistical procedures. The aim of this comparison is to identify methods suitable for these specific gripper-design
problems rather than a universally superior optimizer.

\subsection{Comparative Synthesis and Research Gap}
\label{subsec:research_gap}

The reviewed literature reveals three related research gaps. First, grasp-planning studies generally optimize grasp poses for predefined end-effectors
rather than the gripper configuration itself. Second, existing multi-suction studies generally assume predefined cup layouts. For parallel-jaw
design, Reichert \cite{reichert2024greiferoptimierung} provides the spline-based formulation used in this thesis but considers only
two fixed evaluation configurations. Third, existing algorithm comparisons do not determine which optimization methods are best
suited to the two geometric gripper-design problems considered in this thesis.

For the multi-suction gripper, the unresolved problem is to optimize the positions of the individual cups to maximize grasp coverage on every
workpiece and the area formed by the cup arrangement. For the two-finger gripper, this thesis continues Reichert's framework by optimizing a rigid
spline-based jaw contour across a predefined grasp points. Each grasp point contributes an individual geometric conformity objective. Both
formulations are expensive multi-objective black-box problems, but their different representations, numbers of objectives, and evaluation procedures
may favour different optimization methods.

Within the literature reviewed for this thesis, no controlled comparison of evolutionary and Bayesian multi-objective algorithms was identified for
these two gripper-design formulations.

\subsection{Research Positioning and Contribution}
\label{subsec:research_positioning}

This thesis addresses these gaps by treating gripper geometry as the primary optimization variable. The multi-suction formulation optimizes the
spatial arrangement of individual suction cups, whereas the two-finger formulation optimizes spline-based contours for rigid parallel jaws. Candidate
designs are assessed using problem-specific geometric and collision-aware evaluations.

A modular framework separates geometric evaluation from optimization through a common interface. This allows the applicable algorithms to operate on
the same problem formulation, constraints, and evaluation conditions. Selected evolutionary and Bayesian multi-objective algorithms are compared under
defined evaluation budgets using repeated independent runs, convergence histories, multi-objective performance indicators, and statistical tests.
Evaluation counts and optimizer overhead are reported separately where computational cost is relevant.

The contribution of this thesis is therefore a problem-specific comparison rather than a claim of universal algorithm superiority. The results provide
evidence for selecting optimization methods for the suction-layout and rigid-jaw-contour formulations considered here.